% Introduction to the EU - Week 9 Reading Summary
% Compile: pdflatex 8b-The-Judiciary-reading-summary.tex

\documentclass[11pt,a4paper]{article}
\usepackage[utf8]{inputenc}
\usepackage[T1]{fontenc}
\usepackage[margin=2.5cm]{geometry}
\usepackage{booktabs}
\usepackage{enumitem}
\setlist{nosep,leftmargin=*}
\usepackage{parskip}
\usepackage{microtype}
\usepackage{hyperref}
\usepackage{xcolor}
\definecolor{eublue}{RGB}{0,51,153}

\hypersetup{
  colorlinks=true,
  linkcolor=eublue,
  urlcolor=eublue,
  pdftitle={8b The Judiciary Reading Summary},
  pdfauthor={Introduction to the EU, UCM}
}

\begin{document}

\title{Andrew Duff: \textit{The Judiciary}\\Reading Summary}
\author{Introduction to the European Union\\Universidad Complutense de Madrid\\Week 9}
\date{}
\maketitle
\vspace{-1em}
\hrule
\vspace{1em}

\section*{Bibliographic Information}
\begin{itemize}
\item \textbf{Author:} Andrew Duff
\item \textbf{Chapter:} The Judiciary (Chapter 5)
\item \textbf{Book:} Constitutional Change in the European Union
\item \textbf{Publisher/Year:} The Author(s), 2022
\item \textbf{DOI:} \url{https://doi.org/10.1007/978-3-031-10665-1_5}
\item \textbf{Source file:} 8b The Judiciary.pdf
\end{itemize}

\section*{Thesis}
Duff argues that the Court of Justice is central to European integration and should evolve toward a federal supreme-court role, but this requires stronger rule-of-law compliance, broader citizen access to justice, and treaty reform that removes jurisdictional blind spots.

\section*{Key Arguments}
\subsection*{I. Constitutional Function Of The Court}
\begin{itemize}
\item The Court's treaty mission is to ensure that law is observed in treaty interpretation/application.
\item The Court works through direct actions, infringement control, annulment review, and preliminary rulings.
\item As Commission politics becomes more overt, legal guardianship shifts further to judicial institutions.
\end{itemize}

\subsection*{II. Primacy, Direct Effect, And Multi-Level Tension}
\begin{itemize}
\item EU law primacy is treated as essential to legal uniformity.
\item The EU is not a classic federation, so enforcement often relies on iterative judicial dialogue with national courts.
\item German constitutional jurisprudence (e.g., Maastricht/Lisbon context, ECB litigation) illustrates enduring limits claims by national courts.
\end{itemize}

\subsection*{III. Charter And Rights Architecture}
\begin{itemize}
\item The Charter of Fundamental Rights is a milestone but underused relative to potential.
\item Duff criticizes the impasse over EU accession to the ECHR framework.
\item He proposes Charter modernization (including environmental rights relevance) and broader rights-oriented jurisprudence.
\end{itemize}

\subsection*{IV. Access, Structure, And Scope Reform}
\begin{itemize}
\item Individual standing remains difficult; access to justice should become more practical.
\item Court overload suggests structural reform options (regional layers/appellate pathways).
\item Jurisdictional exclusions (especially around quasi-intergovernmental instruments and sensitive policy fields) reduce federal-court capacity.
\end{itemize}

\section*{Conceptual Tensions}
\begin{itemize}
\item \textbf{Uniformity vs legitimacy:} stronger integration claims can trigger sovereignty backlash.
\item \textbf{Judicial authority vs democratic anchoring:} federal-style adjudication needs parallel political legitimation.
\item \textbf{Rights expansion vs institutional rivalry:} ECJ-Strasbourg coordination remains unresolved.
\end{itemize}

\section*{Relevance To Course}
\begin{itemize}
\item Connects institutional design to real rule-of-law enforcement limits.
\item Bridges legal doctrine (primacy/direct effect) with political conflict among member states.
\item Clarifies why court reform debates are constitutional, not only technical.
\end{itemize}

\section*{Discussion Questions}
\begin{enumerate}
\item Can the Court become a de facto federal supreme court without explicit federal constitutional redesign?
\item Does broadening standing improve democratic legitimacy, or mainly increase judicialization?
\item Should primacy be constitutionalized explicitly in future treaty reform despite likely political resistance?
\end{enumerate}

\vspace{1em}
\noindent\footnotesize\textit{Introduction to the European Union, Universidad Complutense de Madrid.}

\end{document}
